\documentclass{article}
\usepackage[UTF8]{ctex}
\usepackage{amsmath,amssymb,amsfonts}
\usepackage{enumitem}
\usepackage{geometry}
\geometry{a4paper, margin=1in}
% === 错题本排版优化（不参与题目标识） ===
\setlength{\parindent}{0pt}
\setlength{\parskip}{0.35em plus 0.10em minus 0.08em}
\setlength{\emergencystretch}{3em}
\setlist{itemsep=0.25em, topsep=0.40em, parsep=0pt}
\allowdisplaybreaks[2]
% === 错题本排版优化结束 ===
\begin{document}
\section*{中国海洋大学2026年高等代数真题}
\section*{117．中国海洋大学2026年高等代数试卷（回忆版）}
（注．题目来自考生回忆，部分表述或数据可能有误）

\subsection{题目1}
1．（可能有误）若 \(f(x)=x^{3}+(1-\mathrm{i}) x^{2}+2 \mathrm{i} x-10-2 \mathrm{i}\) 有实数根，求 \(f(x)\) 的所有实数根．

\subsection{题目2}
2．已知 \(x=y z\) ，求行列式

\[
D_{n}=\left|\begin{array}{cccc}
1+x & z & & \\
y & 1+x & \ddots & \\
& & & z \\
& & y & 1+x
\end{array}\right|
\]

\subsection{题目3}
3．（可能有误）设 \(\alpha_{1}, \alpha_{2}, \alpha_{3}\) 是 \(\mathbb{R}^{3}\) 的一组基，且

\[
\beta_{1}=2 \alpha_{1}+(2+k) \alpha_{3}, \beta_{2}=2 \alpha_{1}, \beta_{3}=\alpha_{1}+(k+1) \alpha_{3} .
\]

（1）证明：\(\beta_{1}, \beta_{2}, \beta_{3}\) 也是 \(\mathbb{R}^{3}\) 的一组基。\\
（2）求 \(k\) 的值，使得存在非零向量 \(\xi\) ，满足 \(\xi\) 在 \(\alpha_{1}, \alpha_{2}, \alpha_{3}\) 与 \(\beta_{1}, \beta_{2}, \beta_{3}\) 下有相同的坐标，求出所有非零的向量 \(\xi\) 。

\subsection{题目4}
4．已知 \(A=\left(\begin{array}{cc}O & A_{1} \\ A_{2} & O\end{array}\right)\) ，其中 \(A_{1}=\operatorname{diag}\left\{1, \frac{1}{2}, \cdots, \frac{1}{n-1}\right\}, A_{2}=\frac{1}{n}\) ．求 \(\sum_{i=1}^{n} \sum_{j=1}^{n} A_{i j}\) ．

\subsection{题目5}
5．已知 \(f(X)=X^{\mathrm{T}} A X\) 是关于 \(x_{1}, x_{2}, \cdots, x_{n}\) 的一个 \(n\) 元实二次型，且实对称矩阵 \(A\) 的所有特征值为\(\lambda_{1}, \lambda_{2}, \cdots, \lambda_{n}\) ，证明：\(f(X)\) 在条件 \(x_{1}^{2}+x_{2}^{2}+\cdots+x_{n}^{2}=1\) 下的最大值为 \(\max _{1 \leq i \leq n} \lambda_{i}\) ．

\subsection{题目6}
6．已知矩阵 \(A=\left(\begin{array}{ccc}0 & 0 & 0 \\ 2025 & 24 & 0 \\ 23 & 0 & 0\end{array}\right)\) ，证明：矩阵方程 \(X^{2}=A\) 无解，其中 \(X\) 为 3 阶矩阵。

\subsection{题目7}
7．已知 \(J\) 是所有元素为 1 的 \(n\) 阶矩阵，\(f(x)=a+b x\) 是有理系数多项式，且 \(A=f(J)\) 。

\begin{itemize}
  \item[（1）] 求 \(J\) 的所有特征值及特征向量．
  \item[（2）] 求 \(A\) 的所有特征子空间。
  \item[（3）] 判断 \(A\) 是否可以对角化？即是否存在可逆矩阵 \(P\) ，使得 \(P^{-1} A P\) 为对角矩阵．
\end{itemize}

\subsection{题目8}
8．已知 \(X\) 是 \(n\) 阶方阵，若存在正整数 \(k\) 使得 \(X^{k}=O\) ，则称 \(X\) 是幂零矩阵。证明：

\begin{itemize}
  \item[（1）] 若 \(A, B\) 是可交换的幂零矩阵，则 \(A+B\) 也是幂零矩阵。
  \item[（2）] 已知 \(A, B\) 是 \(n\) 阶方阵，若 \(A B\) 是幂零矩阵，则 \(B A\) 也是幂零矩阵．
  \item[（3）] （可能有误）若 \(A\) 是幂零矩阵，证明：不存在 \(n\) 阶方阵 \(B\) ，使得 \(A=B^{2}\) 。
  \item[（4）] 若 \(A\) 是对称矩阵，则 \(A\) 是幂零矩阵的充要条件为 \(A=O\) 。
\end{itemize}

\subsection{题目9}
9．（可能有误）设 \(\varphi\) 是 \(n\) 维欧氏空间 \(V\) 上的线性变换，\(U\) 是 \(V\) 的子空间，且 \(U\) 是 \(\varphi\) 的不变子空间，证明：\(U^{\perp}\) 也是 \(\varphi\) 的不变子空间。
\end{document}
